\section{Runtime Considerations}
\label{sec:runtime}

Sparsification's traditional justification is speed, so we report the runtime picture honestly; it is less favorable than commonly assumed.

\paragraph{Clustering time does not track edge count.}
At realized retention near one half (nominal $\alpha = 0.8$, with-replacement sampler), Leiden on the sparsified graph was \emph{slower} than on the original on four of six large networks in our measurements (speedup $0.58$--$0.88$), with genuine speedups only on the two largest (\texttt{cit-Patents}, \texttt{wiki-topcats}).
The cause is structural: sparsification fragments the graph and inflates the number of detected communities (by factors of $20$--$40$ in the ground-truth experiments of Section~\ref{sec:artifact-granularity}), and Leiden's per-community bookkeeping grows accordingly.
At aggressive retention the fragmentation regime dominates and runtime \emph{increases} despite the smaller edge set.
Speedup claims should therefore be read jointly with fragmentation diagnostics (component counts, largest-component fraction, community counts), and the fragmented regime treated as out of scope for quality claims of any kind.

\paragraph{Pipeline accounting.}
DSpar scoring itself is genuinely cheap---$4$--$8\%$ of end-to-end pipeline time in our measurements, consistent with its $O(m)$ design---so the pipeline cost is dominated by clustering, and the seeded-refinement pipeline of Section~\ref{sec:boundary} costs $1.55$--$1.90\times$ a single baseline run.
This is exactly why Artifact~III matters: the honest comparison for any such pipeline is against a baseline granted the same budget.

\paragraph{Spectral baselines.}
We do not report effective-resistance sparsification comparisons.
In auditing the earlier version of this work we found its spectral branch had silently returned the original graph unmodified (identical edge counts and clustering times), invalidating the previously reported comparison; a corrected implementation is left to future work, with the protocol of Section~\ref{sec:protocol} defining what such a comparison must control for.
